
\documentclass{article}

\usepackage[margin=0.75in]{geometry}
\usepackage{amsmath,amssymb,mathtools}
\usepackage{hyperref}

\title{Structured Representation Learning for Low-Resource Speech Coding}
\author{Ramakrishna Sen}
\date{\today}

\begin{document}

\maketitle

\section{Problem Statement}

We investigate whether audio admits a compact, structured,
and reconstructible representation whose constituent information
can be separated, organized, related, and efficiently encoded.

Unlike conventional neural codecs, which primarily optimize a
latent representation for rate--distortion performance, our
objective is to investigate the structure of the representation
itself. Specifically, we seek to determine what constituent
information is required to describe speech, how this information
can be organized into distinct representations, and how these
representations can jointly reconstruct the original signal.

We propose representing speech through a finite collection of
\emph{feature slots}. Each slot acts as a container for a
particular class of information extracted from the input.
The contents and relationships of these slots collectively
constitute the representation of the audio.

Importantly, the constituent features need not be statistically
independent or mutually orthogonal. Instead, the representation
should minimize unnecessary redundancy while retaining
sufficient information for reconstruction.

The long-term objective is to discover these constituent
features automatically. However, our initial investigation
uses known acoustic properties of human speech to establish
the feasibility of the proposed representation.

\section{Feature-Based Speech Representation}

Consider an audio signal divided into $T$ frames,

\begin{equation}
    X = \{\mathbf{x}_1,\mathbf{x}_2,\ldots,\mathbf{x}_T\},
    \qquad \mathbf{x}_t\in\mathbb{R}^{N}.
\end{equation}

We introduce $K$ feature slots,

\begin{equation}
    \mathcal{C}
    =
    \{\mathcal{C}_1,\mathcal{C}_2,\ldots,\mathcal{C}_K\},
\end{equation}

where each slot contains information extracted from the
audio through a corresponding analysis operator.

For the initial supervised formulation, we define

\begin{equation}
    \boxed{
    \mathbf{z}_{k,t}
    =
    E_k(\mathbf{x}_{\leq t};\theta_k)
    }
\end{equation}

where $E_k$ is a learned feature extractor,
$\theta_k$ denotes its parameters, and
$\mathbf{z}_{k,t}\in\mathcal{Z}_k$ is the representation
associated with slot $k$ at time $t$.

Each slot therefore contains a time-indexed representation,

\begin{equation}
    \mathcal{C}_k(X)
    =
    \{\mathbf{z}_{k,1},\ldots,\mathbf{z}_{k,T}\}.
\end{equation}

The complete structured representation is

\begin{equation}
    \boxed{
    \mathcal{Z}(X)
    =
    \{\mathcal{C}_1(X),\ldots,\mathcal{C}_K(X)\}.
    }
\end{equation}

The individual latent spaces $\mathcal{Z}_k$ need not have
the same dimensionality or represent the same type of
acoustic information.

\subsection{Initial Feature-Slot Definitions}

For the first experiment, we consider the following
known acoustic properties of human speech:

\begin{enumerate}
    \item \textbf{Pitch and voicing:}
    Fundamental-frequency contours and the presence
    or absence of periodic excitation.

    \item \textbf{Spectral envelope:}
    Slowly varying spectral characteristics associated
    with vocal-tract filtering and formant structure.

    \item \textbf{Aperiodicity:}
    Noise-like excitation components and departures
    from harmonic periodicity.

    \item \textbf{Transient structure:}
    Rapidly varying acoustic events associated with
    articulation, consonant bursts, and temporal changes.
\end{enumerate}

These are not assumed to be independent components.
For example, pitch and voicing are closely related,
while spectral-envelope characteristics contribute to
both phonetic identity and perceived timbre.

An additional residual representation may be introduced
to capture information inadequately represented by the
initial feature slots.

The purpose of this initial selection is not to claim
that these properties constitute an optimal decomposition
of speech, but to establish a measurable starting point.

\section{Feature Extraction and Reconstruction}

Each feature extractor is initially trained to estimate
its corresponding acoustic parameters.

Given reference feature representations
$\mathbf{z}_{k,t}^{*}$ obtained from established
signal-processing methods, a preliminary training
objective is

\begin{equation}
    \mathcal{L}_{\mathrm{feat}}
    =
    \sum_{k=1}^{K}
    \alpha_k
    \mathbb{E}_t
    \left[
    d_k\left(
    E_k(\mathbf{x}_{\leq t}),
    \mathbf{z}_{k,t}^{*}
    \right)
    \right],
\end{equation}

where $d_k$ is an appropriate feature-specific
distortion measure.

The extracted representations are combined using a
shared synthesis operator $S$:

\begin{equation}
    \boxed{
    \hat{\mathbf{x}}_t
    =
    S(\mathbf{z}_{1,\leq t},
      \ldots,
      \mathbf{z}_{K,\leq t}).
    }
\end{equation}

The first objective is to establish whether the
selected representations contain sufficient information
for reconstructing speech before introducing quantization.

This can initially be evaluated using classical
analysis--synthesis systems, such as WORLD
\cite{world}, which represents speech through
fundamental frequency, spectral envelope,
and aperiodicity.

Once individual extraction performance has been
established, the extractors and synthesizer may be
jointly fine-tuned using a reconstruction objective,

\begin{equation}
    \mathcal{L}_{\mathrm{rec}}
    =
    \mathbb{E}
    \left[
    D\left(
    \mathbf{x},
    \hat{\mathbf{x}}
    \right)
    \right],
\end{equation}

where $D$ may combine waveform and perceptually
motivated spectral distortion measures.

This distinction is important: accurate estimation
of individual acoustic parameters does not necessarily
guarantee accurate joint waveform reconstruction.

\section{Quantization and Transmission}

Having established a reconstructible representation,
we investigate how efficiently its constituent
information can be transmitted.

Each feature representation is quantized independently
in the initial experiment:

\begin{equation}
    \tilde{\mathbf{z}}_{k,t}
    =
    Q_k(\mathbf{z}_{k,t}).
\end{equation}

The decoder reconstructs the audio using only the
received quantized representations:

\begin{equation}
    \boxed{
    \hat{\mathbf{x}}_t
    =
    S(\tilde{\mathbf{z}}_{1,\leq t},
      \ldots,
      \tilde{\mathbf{z}}_{K,\leq t}).
    }
\end{equation}

Because the initial slots have fixed, predefined
meanings, their identities need not be transmitted
explicitly. The bitstream instead contains the
quantized feature parameters and any necessary
timing, scaling, or activity information.

For a fixed-rate codec, the total transmitted bits
over a coding interval must satisfy

\begin{equation}
    \boxed{
    B_{\mathrm{total}}
    =
    \sum_{k=1}^{K} B_k
    +
    B_{\mathrm{overhead}}
    \leq B_{\max}.
    }
\end{equation}

We therefore investigate the constrained
rate--distortion problem

\begin{equation}
    \boxed{
    \begin{aligned}
    \min_{\{Q_k\},S}\quad&
    \mathbb{E}[D(\mathbf{x},\hat{\mathbf{x}})]
    \\
    \text{subject to}\quad&
    B_{\mathrm{total}}\leq B_{\max}.
    \end{aligned}
    }
\end{equation}

Different slots may require different quantization
strategies and bit allocations. The objective is
to establish the marginal reconstruction benefit
of allocating bits to each feature family.

If the initial feature representation is insufficient,
an additional learned residual slot may be introduced
to encode information missing from the base reconstruction.

\section{Extension to Learned Slot Discovery}

The preceding formulation assumes that the contents
of each slot are known beforehand.

Our longer-term objective is to remove this assumption
and allow the architecture to discover useful
constituent representations automatically.

We propose first extracting a collection of learned
feature tokens from each audio frame,

\begin{equation}
    \mathcal{H}_t
    =
    \{\mathbf{h}_{t,1},\ldots,\mathbf{h}_{t,L}\}.
\end{equation}

A learned assignment mechanism then determines
the association between each token and the available
slots:

\begin{equation}
    \boxed{
    a_{t,n,k}
    =
    P(\mathcal{C}_k\mid\mathbf{h}_{t,n}),
    \qquad
    \sum_{k=1}^{K}a_{t,n,k}=1.
    }
\end{equation}

The representation of slot $k$ is obtained by
aggregating its assigned information,

\begin{equation}
    \mathbf{z}_{k,t}
    =
    \operatorname{Agg}_k
    \left(
    \{a_{t,n,k}\mathbf{h}_{t,n}\}_{n=1}^{L}
    \right).
\end{equation}

This formulation is related to slot-based
representation learning, including Slot Attention
\cite{slotattention}, but differs in its intended
application: rather than identifying discrete
objects or sound sources, we seek constituent
acoustic information suitable for reconstruction
and efficient coding.

Learned slot assignments may require additional
transmission information. Their coding cost must
therefore be included when evaluating the representation.

\section{Functional Separation and Statistical Dependencies}

A major challenge is preventing different slots
from learning redundant representations.

However, conventional orthogonality between
dictionary matrices,

\begin{equation}
    \mathbf{D}_i^\top\mathbf{D}_j=\mathbf{0},
\end{equation}

may be unnecessarily restrictive because
different speech properties are neither physically
nor statistically independent.

Simon et al.~\cite{simon2026} propose an alternative
principle based on the functional effects of
latent variables on a generative mapping.

For a differentiable synthesis function
$\hat{\mathbf{x}}=S(\mathbf{z})$, define

\begin{equation}
    \mathbf{J}_{S,k}
    =
    \frac{\partial S}{\partial\mathbf{z}_k}.
\end{equation}

A possible group-wise functional orthogonality
regularizer is

\begin{equation}
    \boxed{
    \mathcal{L}_{\mathrm{func}}
    =
    \sum_{i<j}
    \left\|
    \mathbf{J}_{S,i}^{\top}
    \mathbf{J}_{S,j}
    \right\|_F^2.
    }
\end{equation}

This encourages changes in different latent groups
to produce complementary local changes in the
reconstructed waveform.

It is an adaptation of the functional-orthogonality
principle, not a direct application of Simon et al.'s
identifiability theorem, which requires additional
assumptions on the generative mapping and latent domain.

Functional separation also does not imply statistical
independence. Consequently, dependencies between
slot representations may still be exploited for coding.

For example, a predictive model could estimate

\begin{equation}
    \hat{\mathbf{z}}_{j,t}
    =
    g_j\left(
    \{\tilde{\mathbf{z}}_{i,\leq t}:
    i\in\operatorname{Pa}(j)\}
    \right),
\end{equation}

where the parent representations are already available
to the decoder.

Only the prediction innovation would then be encoded:

\begin{equation}
    \boxed{
    \begin{aligned}
    \mathbf{e}_{j,t}
    &=
    \mathbf{z}_{j,t}
    -
    \hat{\mathbf{z}}_{j,t},
    \\
    \tilde{\mathbf{z}}_{j,t}
    &=
    \hat{\mathbf{z}}_{j,t}
    +
    Q_j(\mathbf{e}_{j,t}).
    \end{aligned}
    }
\end{equation}

Such dependencies may eventually be represented
using a sparse directed graph. Graph learning and
functional orthogonality are therefore considered
subsequent extensions rather than prerequisites
for the initial representation experiment.

\section{Proposed Research Methodology}

The investigation is organized into four stages.

\begin{enumerate}
    \item \textbf{Known-property representation:}
    Extract pitch, voicing, spectral-envelope,
    and aperiodicity information using established
    analytical methods. Evaluate reconstruction
    before quantization.

    \item \textbf{Learned feature extraction and coding:}
    Train lightweight feature-specific networks
    and a shared synthesizer. Quantize the
    representations and evaluate reconstruction
    at matched bitrates and computational costs.

    \item \textbf{Unsupervised slot discovery:}
    Replace predefined feature assignments with
    a learned routing mechanism. Investigate
    representation sufficiency, slot utilization,
    stability, and redundancy.

    \item \textbf{Functional and statistical structure:}
    Evaluate functional orthogonality constraints
    and subsequently investigate predictive
    relationships between slots for reducing
    transmitted information.
\end{enumerate}

Throughout these experiments, three properties
remain central: reconstruction fidelity,
representation economy, and computational efficiency.

The ultimate research question is whether a
structured representation of audio, built from
learned constituent information and its relationships,
can provide a better rate--distortion--complexity
tradeoff than an unconstrained latent representation.

The objective is not to assume that disentanglement
or orthogonality improves compression, but to
establish whether organizing audio information
into such representations provides a measurable
advantage.

\begin{thebibliography}{9}

\bibitem{slotattention}
F. Locatello et al.,
``Object-Centric Learning with Slot Attention,''
\emph{Advances in Neural Information Processing Systems},
2020.

\bibitem{world}
M. Morise, F. Yokomori, and K. Ozawa,
``WORLD: A Vocoder-Based High-Quality Speech Synthesis
System for Real-Time Applications,''
\emph{IEICE Transactions on Information and Systems},
2016.

\bibitem{simon2026}
M. C. Simon, P. Frossard, and C. De Vleeschouwer,
``Unsupervised Disentanglement Without Compromises:
How Functional Orthogonality Enforces Identifiability,''
\emph{Proceedings of the 43rd International Conference
on Machine Learning}, 2026.

\end{thebibliography}

\end{document}
